\chapter{The relative Picard functor}
\label{chap:Picard_Functor}

This chapter sets up Phase~C step~2: the \emph{relative Picard functor} of a curve $C$ over $\Spec k$, packaged as a contravariant functor on schemes, and the (deferred) representability theorem whose closure delivers the Picard scheme. It is the natural successor to Chapter~\ref{chap:Picard_LineBundle}.

The chapter introduces no protected declarations: every definition and lemma here is auxiliary, supporting the construction of $\Jac(C)$ in Chapter~\ref{chap:Jacobian}.

\section{Definition}

% NOTE (iter-005 review): The Lean realization takes the source category as
% (Over (Spec (CommRingCat.of k)))ᵒᵖ rather than the bare Schemeᵒᵖ informally
% written below — a generic S : Scheme has no structure morphism to Spec k, so
% the fiber product C ×_k S is not canonically defined without the over-Spec k
% bundling. This matches the chapter's mathematical setup "Sch_k = Sch / Spec k"
% and the convention already used in Jacobian.lean. See task_results/Picard_Functor.lean.md.
% NOTE (iter-005 review): The current definition rests on the global-sections
% approximation of LineBundle from Chapter~\ref{chap:Picard_LineBundle} — see
% the forward-compatibility section below.
\begin{definition}[Relative Picard functor]
  \label{def:Pic_functor}
  \lean{AlgebraicGeometry.Scheme.PicardFunctor}
  \leanok
  Let $C$ be a scheme over $\Spec k$. The \emph{relative Picard functor} of $C$ over $k$ is the contravariant functor
  \[
    \PicardFunctor\, C \;\colon\; \mathrm{Sch}^{\mathrm{op}} \;\longrightarrow\; \mathrm{Set},\qquad
    S \;\longmapsto\; \Pic(C \times_k S) \,\big/\, p_S^* \Pic(S),
  \]
  where $p_S \colon C \times_k S \to S$ is the projection and the quotient is by the subgroup $p_S^*\Pic(S) \subseteq \Pic(C \times_k S)$ pulled back from $S$. Functoriality in $S$ is via base change of the fiber product and the pull-back homomorphism on Picard groups (Theorem~\ref{thm:Scheme_Pic_pullback}).
\end{definition}

\begin{remark}
  The presheaf $\PicardFunctor$ above is sometimes called the \emph{absolute} Picard presheaf; the standard \emph{relative Picard functor} of $C/k$ is its étale-sheafification, denoted $\underline{\Pic}_{C/k}$. For a smooth proper geometrically connected curve $C$ over a field $k$, étale sheafification is automatic on schemes that are themselves Picard sheaves over $S$ (the obstruction is the Brauer group), so this chapter records the presheaf definition; the étale sheafification is deferred to a subsequent iteration (Phase~C step~3).
\end{remark}

\section{Representability}

% NOTE (iter-005 review): Statement is formalized (Lean declaration exists,
% file compiles), proof is intentionally a sorry per the forward-compatibility
% section below — closing it on the global-sections-approximate LineBundle
% would silently assert representability of the wrong functor. No \leanok on
% the proof block until Phase B/C deepens LineBundle.
\begin{theorem}[Representability of the Picard functor — Grothendieck FGA, Mumford \emph{Abelian Varieties} §III.13]
  \label{thm:Pic_representable}
  \lean{AlgebraicGeometry.Scheme.PicardFunctor.representable}
  \leanok
  \uses{def:Pic_functor}
  Let $C$ be a smooth proper geometrically irreducible curve over a field $k$. Then $\PicardFunctor\, C$ (after étale sheafification) is representable by a smooth, locally finite-type group scheme $\Pic_{C/k}$ over $k$. Its connected component of the identity $\Pic^0_{C/k}$ is a smooth proper geometrically irreducible group scheme over $k$ of dimension $g = g(C)$, hence an abelian variety.
\end{theorem}

\begin{proof}
  \uses{def:Pic_functor, def:Scheme_LineBundle, thm:Scheme_Pic_pullback}
  See \emph{Fundamental Algebraic Geometry: Grothendieck's FGA Explained} (Fantechi--Göttsche--Illusie--Kleiman--Nitsure--Vistoli, 2005), Chapter~9, and Mumford, \emph{Abelian Varieties}, Chapter~III §13. The proof goes via Hilbert schemes and quotient stacks, both of which currently exceed the Mathlib substrate. The Lean side leaves this as a deferred sorry; downstream consumers (\ref{def:Jacobian} and the four \lstinline{Jacobian} instances of Chapter~\ref{chap:Jacobian}) record their dependence on it.
\end{proof}

\section{Use in the project}

Theorem~\ref{thm:Pic_representable} is the bundled blocker for the four \lstinline{Jacobian.lean} sorries:

\begin{itemize}
  \item Definition~\ref{def:Jacobian} ($\Jac(C) = \Pic^0_{C/k}$) takes the connected component of identity of the representing scheme.
  \item Theorem~\ref{thm:Jacobian_grpObj} (group-object structure) is the abelian-variety structure on $\Pic^0_{C/k}$.
  \item Theorem~\ref{thm:Jacobian_smooth_genus} (smooth of relative dimension $g$) is the dimension formula in the representability statement.
  \item Theorem~\ref{thm:Jacobian_proper} (properness) and Theorem~\ref{thm:Jacobian_geomIrred} (geometric irreducibility) are direct consequences of the representability conclusion.
\end{itemize}

\section{Mathlib gap}

Mathlib \texttt{b80f227} has the constituent infrastructure for the \emph{definition} (Definition~\ref{def:Pic_functor}):
\begin{itemize}
  \item Fiber products of schemes (\lstinline{CategoryTheory.Limits.HasPullback} for the category of schemes).
  \item \lstinline{AlgebraicGeometry.Scheme.LineBundle} (this project, Chapter~\ref{chap:Picard_LineBundle}).
  \item \lstinline{AlgebraicGeometry.Scheme.Pic.pullback} (this project, Chapter~\ref{chap:Picard_LineBundle}).
  \item \lstinline{QuotientGroup.mk}, \lstinline{QuotientAddGroup.mk} for group quotients.
  \item \lstinline{CategoryTheory.Functor} for assembling the contravariant functor.
\end{itemize}

What is \emph{missing} and supplied by this chapter (as a deferred sorry):
\begin{itemize}
  \item The actual representability theorem (Theorem~\ref{thm:Pic_representable}). FGA-level, kept as a sorry; all four \lstinline{Jacobian.lean} sorries reduce to it once it is closed.
\end{itemize}

\subsection{Forward-compatibility note (\texttt{LineBundle} approximation)}

Chapter~\ref{chap:Picard_LineBundle} defines \lstinline{LineBundle X} as the global-sections approximation \lstinline{CommRing.Pic Γ(X, ⊤)}, correct on affine schemes (Stacks 0AGS) but a strict subgroup of the true $\Pic(X)$ on non-affine schemes. The \emph{definition} of \lstinline{PicardFunctor C} therefore inherits this approximation: the values $\PicardFunctor(S)$ for non-affine $S$ are smaller subgroups of the true relative Picard group. Consequently the representability theorem stated above will be honest only after \lstinline{LineBundle} is refined to the bespoke "invertible quasi-coherent sheaf" definition (gated on Mathlib's \lstinline{MonoidalCategory X.Modules}, currently absent). Until that refinement lands, the \lstinline{representable} sorry of this chapter must remain unattacked: closing it on the approximate \lstinline{LineBundle} would silently assert representability of the wrong functor.

\subsection{Discovery objectives for this iteration (iter-005 prover)}

The prover assigned to this chapter should:
\begin{enumerate}
  \item Fill the body of \lstinline{AlgebraicGeometry.Scheme.PicardFunctor C} as the contravariant functor $S \mapsto \Pic(C \times_k S) / p_S^*\Pic(S)$. Use:
    \begin{itemize}
      \item \lstinline{Limits.pullback} for the fiber product $C \times_k S$ in the category of schemes;
      \item \lstinline{AlgebraicGeometry.Scheme.LineBundle} (now closed in Chapter~\ref{chap:Picard_LineBundle}) for $\Pic(C \times_k S)$ and $\Pic(S)$;
      \item \lstinline{AlgebraicGeometry.Scheme.Pic.pullback} for the homomorphism $\Pic(S) \to \Pic(C \times_k S)$;
      \item \lstinline{QuotientGroup.mk} for the quotient $\Pic(C \times_k S) / p_S^*\Pic(S)$;
      \item \lstinline{QuotientGroup.lift} for the induced map on the quotients (functoriality of $\PicardFunctor$ in $S$).
    \end{itemize}
    Document each constituent in the task result.
  \item \textbf{Do not} attempt to fill the \lstinline{representable} sorry. It is FGA-level and intentionally deferred. Even if a "trivial" closure is found (e.g.\ via the global-sections approximation), it would assert representability of the wrong functor and is therefore a forbidden shortcut. Keep \lstinline{sorry} in place.
  \item Add functoriality lemmas (\lstinline{PicardFunctor.map_id}, \lstinline{PicardFunctor.map_comp}) only if they are not already automatic from \lstinline{Functor.map_id} / \lstinline{Functor.map_comp}; in most cases the bundled \lstinline{Functor} structure provides them for free. Decompose into private lemmas if a one-shot definition exceeds $\sim 60$ lines.
  \item If a Mathlib gap blocks progress on the definition (e.g.\ "the pull-back of fiber products of schemes is not exposed in this exact form"), document the missing declaration in the task result with name and intended type — \emph{do not} fill the sorry with an ad-hoc placeholder.
\end{enumerate}
